\documentclass[../ai_in_finance]{subfiles}
\usepackage[margin=1in]{geometry}
\usepackage{tikz}
\usetikzlibrary{positioning}
\usepackage{amsmath}
\usepackage{amssymb}
\usepackage{booktabs}
\usepackage{array}
\newcolumntype{L}[1]{>{\raggedright\arraybackslash}p{#1}}
\usepackage{url}
\usepackage{draftwatermark}

\SetWatermarkText{Wulin.Suo@Queensu.ca}
\SetWatermarkScale{1.4}
\SetWatermarkColor{gray!35}
\SetWatermarkAngle{90}
\SetWatermarkHorCenter{0.94\paperwidth}
\SetWatermarkVerCenter{0.5\paperheight}

\usepackage{makeidx}
\makeindex

\begin{document}

\ifSubfilesClassLoaded{\appendix}{}
\chapter{A Career Map: Roles, Chapters, and Preparation}
\label{app:careers}

This book's stated audience is a reader with quantitative training who wants a career in the financial industry, and the chapters have been building the technical vocabulary that industry runs on. What the chapters do not do, because no single chapter is the right place for it, is describe the industry's actual job landscape: which roles exist, what each one does all day, and which parts of this book matter most for each. This appendix closes that gap. It is a map, not a syllabus; job titles vary considerably across firms (the same work may be called ``quantitative researcher'' at a hedge fund, ``strats'' at an investment bank, and ``data scientist'' at a fintech), so the descriptions below are organized around what the work actually is rather than what any particular firm calls it.\index{careers in financial AI}

One framing is worth internalizing before the table. Broadly, quantitative roles in finance divide into those that \emph{make} risk (research and trading roles, judged on returns), those that \emph{measure and police} risk (risk management, model validation, and compliance roles, judged on rigor and independence), and those that \emph{build the machinery} both sides run on (engineering and data roles, judged on reliability and scale). All three groups hire people with this book's intended background, the mix of statistics, programming, and financial context, but they reward different temperaments, and the interview processes differ accordingly: the first group leans hardest on probability, brainteasers, and market intuition; the second on statistical rigor, regulatory awareness, and the ability to explain a model's limits; the third on software engineering and system design, with the finance tested more lightly.

\begin{table}[htbp]
\centering
\begin{lrbox}{\bookmeasuredbox}
\begin{tabular}{L{3.4cm}L{5.6cm}L{3.4cm}}
\toprule
Role & What the job actually is & Core chapters \\
\midrule
Quantitative researcher (buy side) & Finding and testing predictive signals in market, fundamental, text, and alternative data; building and validating backtests & 2, 5--7, 12, 13, 16, 17 \\
Execution and e-trading quant & Designing the algorithms that execute orders and make markets: impact models, order routing, market-making logic & 3, 13, 14, 19 \\
Market risk quant & Measuring how much the firm can lose to market moves: VaR and ES models, stress testing, backtesting, limits & 5, 7, 8 \\
Credit risk modeler & Estimating who will default and what it costs: PD/LGD/EAD models, scorecards, portfolio credit models, capital & 4, 6, 9, 18, 21 \\
Actuary and insurance data scientist & Pricing, reserving, claims, and catastrophe modeling for insurers: GLMs and their machine-learning challengers inside a filed-rating regime & 5, 6, 10, 11, 18 \\
Model validator (second line) & Independently challenging every model above: testing assumptions, benchmarking, documentation, SR 11-7 discipline & 6--11, 18 \\
Operational-risk and surveillance analyst & Quantifying and detecting operational failures: loss models, anomaly detection, conduct and communications surveillance & 6, 10, 15, 16 \\
Fraud and financial-crime data scientist & Catching fraud, money laundering, and market abuse in transaction and network data, at extreme class imbalance & 6, 14, 15, 17, 18 \\
NLP and alternative-data researcher & Turning filings, calls, news, satellite and web data into investable signals with point-in-time discipline & 13, 16, 17 \\
Fintech and lending ML engineer & Building production credit, robo-advisory, and payments models end to end, with fairness and explainability constraints & 2, 6, 9, 12, 18, 21 \\
AI governance and compliance specialist & Setting and enforcing the rules AI systems must satisfy: model inventories, monitoring, fairness audits, regulatory response & 10, 15, 16, 18 \\
\bottomrule
\end{tabular}
\end{lrbox}
\ifdim\wd\bookmeasuredbox<0.65\linewidth
\setlength{\bookcapwidth}{0.65\linewidth}
\else
\setlength{\bookcapwidth}{\wd\bookmeasuredbox}
\fi
\floatbox[\capbot]{table}[\bookcapwidth]{
\usebox{\bookmeasuredbox}
}{
\caption{A map from common quantitative-finance roles to the chapters that prepare for them}
\label{tab:appendix-careers}
}
\end{table}

\section{Reading the Map}

A few patterns in Table~\ref{tab:appendix-careers} deserve comment. Chapters 2 and 6, Python and machine learning fundamentals, appear implicitly in every row; they are the price of admission everywhere and are omitted from rows only where a role's distinctive preparation lies elsewhere. Chapter 18, explainability and fairness, appears in every role that touches a regulated decision about a person, which is the pattern the chapter itself predicts: the closer a model sits to a consumer outcome, the more its explanation apparatus matters relative to its accuracy. And Chapter 21's capstone appears wherever the honest answer to ``what should I build to demonstrate competence?'' is a complete, end-to-end project rather than a familiarity with any single method; a worked capstone, cleaned data and calibration checks and fairness audit included, is the single most credible artifact a candidate without industry experience can bring to an interview, precisely because its eight-stage template mirrors how the work is actually structured inside a firm.

The risk-and-validation column of the industry deserves a special word, because candidates with strong ML backgrounds systematically underrate it. Model validation and risk roles are often easier to enter than front-office research, teach the firm's full model landscape faster than any other seat, are hired against a regulatory requirement rather than a profit forecast (making them notably more stable across market cycles), and increasingly involve genuinely hard technical work, challenging a deep-learning surveillance model requires understanding it at least as well as its builder did. A reader whose favorite chapters were 8 through 11 and 18 should take that as real information about where they would thrive.

\section{Preparing, Concretely}

Three concrete uses of this book map directly onto how candidates are actually evaluated. First, the worked examples are interview questions: computing a two-asset portfolio variance by hand, walking through a limit-order-book fill, explaining why the 99.9\% operational-loss quantile dwarfs its mean, deriving a Shapley decomposition for a small linear model, and reasoning through Bayes' theorem at a realistic base rate are all standard fare, and every one of them is a worked example somewhere in Chapters 1 through 21. Second, the challenge exercises are portfolio projects: each one pairs a method from a chapter with a real, freely downloadable dataset, which is the scope of take-home assignments many firms now use in place of a second interview round. Third, the Further Reading lists are the canon: Hull for derivatives and risk, Bodie-Kane-Marcus for investments, Chan and Hilpisch for systematic trading, and the original papers behind XGBoost, SHAP, and the transformer are the references interviewers themselves learned from, and recognizing them signals fluency the way no self-description can.

The map's final message is the book's own: every role in Table~\ref{tab:appendix-careers} is being reshaped by the same forces this book has traced chapter by chapter, and the durable advantage is not any single tool but the combination the intended reader has been building all along, enough statistics to be skeptical, enough programming to be self-sufficient, and enough financial context to know what the numbers mean and what happens when they are wrong.

\end{document}
